% !TeX program = xelatex
% Kompilacja:  python build.py      (uzywa silnika Tectonic -> main.pdf)
\documentclass[11pt,a4paper]{report}

% ---------------------------------------------------------------------
%  Silnik i czcionki (dziala pod Tectonic/XeTeX oraz pod pdfLaTeX)
% ---------------------------------------------------------------------
\usepackage{iftex}
\usepackage{amsmath}
\usepackage{mathtools}

\ifPDFTeX
  \usepackage[utf8]{inputenc}
  \usepackage[T1]{fontenc}
  \usepackage{lmodern}
  \usepackage{amssymb}
\else
  % XeTeX/LuaTeX (Tectonic): natywny Unicode + eleganckie czcionki OTF
  \usepackage{fontspec}
  \usepackage{unicode-math}
  \setmainfont{texgyrepagella-regular.otf}[
    BoldFont       = texgyrepagella-bold.otf,
    ItalicFont     = texgyrepagella-italic.otf,
    BoldItalicFont = texgyrepagella-bolditalic.otf ]
  \setsansfont{texgyreheros-regular.otf}[
    BoldFont   = texgyreheros-bold.otf,
    ItalicFont = texgyreheros-italic.otf,
    Scale      = 0.92 ]
  \setmathfont{texgyrepagella-math.otf}
\fi

\usepackage[polish]{babel}
\usepackage[a4paper,margin=2.5cm,headheight=15pt]{geometry}
\usepackage{microtype}
\usepackage{parskip}
\linespread{1.05}

% ---------------------------------------------------------------------
%  Kolory
% ---------------------------------------------------------------------
\usepackage{xcolor}
\definecolor{accent}{HTML}{1F3A5F}   % granatowy
\definecolor{accentB}{HTML}{2E7D32}  % zielony
\definecolor{accentBg}{HTML}{EEF3F8}
\definecolor{warn}{HTML}{B26A00}
\definecolor{warnBg}{HTML}{FBF0E1}
\definecolor{softgray}{HTML}{F4F5F7}

% ---------------------------------------------------------------------
%  Grafika, tabele, listy
% ---------------------------------------------------------------------
\usepackage{graphicx}
\usepackage{booktabs}
\usepackage{array}
\usepackage{enumitem}
\setlist{itemsep=2pt, topsep=3pt, parsep=0pt}
\usepackage{caption}
\captionsetup{font=small, labelfont={sf,bf,color=accent}, textfont=sf}

\usepackage{tikz}
\usetikzlibrary{arrows.meta,calc,positioning,intersections,%
  decorations.pathreplacing,patterns,backgrounds,fit}
\usepackage{pgfplots}
\pgfplotsset{compat=1.18}
\usepgfplotslibrary{groupplots,fillbetween}
\tikzset{>=Stealth}

% ---------------------------------------------------------------------
%  Kolorowe ramki (definicje / intuicje / wzory / uwagi)
% ---------------------------------------------------------------------
\usepackage[most]{tcolorbox}
\newtcolorbox{definicja}[1][]{enhanced, breakable, colback=accentBg,
  colframe=accent, coltitle=white, fonttitle=\sffamily\bfseries,
  title=Definicja, boxrule=0.6pt, arc=2pt, left=8pt, right=8pt,
  top=5pt, bottom=5pt, #1}
\newtcolorbox{intuicja}[1][]{enhanced, breakable, colback=accentB!6,
  colframe=accentB, coltitle=white, fonttitle=\sffamily\bfseries,
  title=Intuicja, boxrule=0.6pt, arc=2pt, left=8pt, right=8pt,
  top=5pt, bottom=5pt, #1}
\newtcolorbox{kluczowe}[1][]{enhanced, breakable, colback=accent!7,
  colframe=accent!55!black, coltitle=white, fonttitle=\sffamily\bfseries,
  title=Kluczowe, boxrule=0.6pt, arc=2pt, left=8pt, right=8pt,
  top=5pt, bottom=5pt, #1}
\newtcolorbox{uwaga}[1][]{enhanced, breakable, colback=warnBg,
  colframe=warn, coltitle=white, fonttitle=\sffamily\bfseries,
  title=Uwaga, boxrule=0.6pt, arc=2pt, left=8pt, right=8pt,
  top=5pt, bottom=5pt, #1}

% ---------------------------------------------------------------------
%  Stylizacja naglowkow
% ---------------------------------------------------------------------
\usepackage{titlesec}
\titleformat{\chapter}[display]
  {\normalfont\sffamily\bfseries\color{accent}}
  {\color{accent!65}\large\MakeUppercase{\chaptertitlename}\ \thechapter}
  {4pt}{\Huge}
\titlespacing*{\chapter}{0pt}{-6pt}{16pt}
\titleformat{\section}
  {\normalfont\sffamily\Large\bfseries\color{accent}}
  {\thesection}{0.6em}{}
  [\vspace{-4pt}{\color{accent!25}\titlerule[1pt]}]
\titleformat{\subsection}
  {\normalfont\sffamily\large\bfseries\color{accent!88!black}}
  {\thesubsection}{0.6em}{}

% ---------------------------------------------------------------------
%  Naglowki i stopki stron
% ---------------------------------------------------------------------
\usepackage{fancyhdr}
\pagestyle{fancy}
\fancyhf{}
\fancyhead[L]{\small\sffamily\color{accent!80}\nouppercase{\leftmark}}
\fancyhead[R]{\small\sffamily\color{accent!80}\thepage}
\renewcommand{\headrulewidth}{0.6pt}
\renewcommand{\footrulewidth}{0pt}
\fancypagestyle{plain}{\fancyhf{}%
  \fancyfoot[C]{\small\sffamily\color{accent!80}\thepage}%
  \renewcommand{\headrulewidth}{0pt}}

% ---------------------------------------------------------------------
%  Odsylacze
% ---------------------------------------------------------------------
\usepackage{hyperref}
\hypersetup{colorlinks=true, linkcolor=accent, urlcolor=accentB,
  citecolor=accentB, bookmarksnumbered=true,
  pdftitle={Zagadnienia egzaminacyjne -- Ekonomia i Finanse}}

\setcounter{tocdepth}{1}
\setcounter{secnumdepth}{2}
\raggedbottom

% =====================================================================
\begin{document}
% =====================================================================

\begin{titlepage}
  \centering
  \vspace*{3.2cm}
  {\sffamily\bfseries\color{accent}\Huge Zagadnienia egzaminacyjne\par}
  \vspace{10pt}
  {\color{accent!25}\rule{0.5\linewidth}{1.2pt}}\par
  \vspace{14pt}
  {\sffamily\LARGE Egzamin doktorski\par}
  \vspace{4pt}
  {\sffamily\Large Ekonomia i~Finanse\par}
  \vfill
  {\sffamily\large Opracowanie i~notatki\par}
  \vspace{6pt}
  {\small \today}
\end{titlepage}

\pagenumbering{roman}
\tableofcontents
\cleardoublepage
\pagenumbering{arabic}

% =====================================================================
\chapter{Ekonomia}
% =====================================================================

% =====================================================================
%  ZAGADNIENIE 1 (Ekonomia)
%  Elastyczność; optimum techniczne/ekonomiczne; użyteczność;
%  maksymalizacja zysku i podaż firmy w konkurencji doskonałej.
% =====================================================================
\section{Elastyczność popytu i podaży; optimum techniczne i ekonomiczne; funkcja użyteczności; maksymalizacja zysku i krzywa podaży przedsiębiorstwa na rynku doskonale konkurencyjnym w krótkim i długim okresie}

\begin{intuicja}[title={Mapa zagadnienia}]
Cztery cegiełki mikroekonomii, które razem opisują \emph{jak reaguje rynek} i \emph{jak decydują jego uczestnicy}:
\begin{enumerate}[leftmargin=1.4em]
  \item \textbf{Elastyczność} — o ile procent zmienia się ilość, gdy cena (lub dochód) zmieni się o 1\%.
  \item \textbf{Optimum producenta} — \emph{techniczne} (najwydajniejsze technologicznie) vs \emph{ekonomiczne} (najbardziej zyskowne).
  \item \textbf{Funkcja użyteczności} — jak konsument wybiera koszyk dóbr przy ograniczonym budżecie.
  \item \textbf{Konkurencja doskonała} — skąd bierze się krzywa podaży firmy oraz co ją różni w krótkim i~długim okresie.
\end{enumerate}
\end{intuicja}

% ---------------------------------------------------------------------
\subsection{Elastyczność popytu i podaży}
% ---------------------------------------------------------------------

\textbf{Elastyczność} to miara wrażliwości jednej wielkości na zmianę drugiej, wyrażona jako iloraz \emph{zmian procentowych}. Dzięki temu jest \emph{bezwymiarowa} (niezależna od jednostek: sztuk, kg, złotych), więc porównujemy nią dowolne dobra.

\begin{definicja}[title={Cenowa elastyczność popytu $E_D$}]
\[
E_D \;=\; \frac{\%\,\Delta Q_D}{\%\,\Delta P}
      \;=\; \frac{dQ}{dP}\cdot\frac{P}{Q}
      \;=\; \frac{d\ln Q}{d\ln P}.
\]
Ponieważ popyt maleje z ceną ($dQ/dP<0$), zwykle $E_D<0$; w praktyce mówimy o wartości bezwzględnej $|E_D|$.
\end{definicja}

\begin{intuicja}
Dla informatyka: elastyczność to \emph{pochodna logarytmiczna} $d\ln Q/d\ln P$ — nachylenie krzywej w skali log–log. Jest lokalna (inna w każdym punkcie) i unormowana, więc nie zależy od tego, czy mierzysz w sztukach, czy w tysiącach sztuk.
\end{intuicja}

\begin{table}[h]
\centering
\small
\begin{tabular}{@{}lll@{}}
\toprule
\textbf{Zakres} & \textbf{Nazwa} & \textbf{Interpretacja} \\
\midrule
$|E_D| = 0$        & doskonale nieelastyczny & ilość w ogóle nie reaguje (pion) \\
$0<|E_D|<1$        & nieelastyczny           & ilość reaguje słabo \\
$|E_D| = 1$        & jednostkowy             & zmiana proporcjonalna \\
$|E_D| > 1$        & elastyczny              & ilość reaguje silnie \\
$|E_D| \to \infty$ & doskonale elastyczny    & ilość dowolna przy danej cenie (poziom) \\
\bottomrule
\end{tabular}
\caption{Klasyfikacja cenowej elastyczności popytu.}
\end{table}

\paragraph{Elastyczność wzdłuż liniowej krzywej popytu.} Wbrew intuicji, prosta krzywa popytu ma \emph{różną} elastyczność w każdym punkcie: od $\infty$ u góry do $0$ przy osi ilości, a dokładnie w środku $|E_D|=1$ (rys.~\ref{fig:elast}).

\begin{figure}[h]
\centering
\begin{tikzpicture}[scale=0.95]
  \draw[->] (0,0) -- (7,0) node[right]{$Q$};
  \draw[->] (0,0) -- (0,6.4) node[above]{$P$};
  \draw[very thick, accent] (0,6) -- (6,0);
  \node[accent] at (5.4,0.95) {$D$};
  % midpoint
  \fill[accent] (3,3) circle (2pt);
  \draw[dashed, gray] (3,0) node[below]{$\tfrac12 Q_{\max}$} -- (3,3) -- (0,3);
  \node[right] at (3,3) {\small $|E_D|=1$};
  % regions (z odnosnikami, poza linia popytu)
  \node[accentB, align=center, font=\small] (elx) at (4.75,5.0) {$|E_D|>1$\\ elastyczny};
  \draw[accentB, ->, thin] (elx) -- (1.35,4.65);
  \node[warn, align=center, font=\small] (inelx) at (1.55,1.15) {$|E_D|<1$\\ nieelastyczny};
  \draw[warn, ->, thin] (inelx) -- (4.5,1.5);
  \node[left] at (0,6) {\small $P_{\max}$};
  \node[below] at (6,0) {\small $Q_{\max}$};
\end{tikzpicture}
\caption{Elastyczność zmienia się wzdłuż liniowej krzywej popytu.}
\label{fig:elast}
\end{figure}

\paragraph{Elastyczność a utarg (przychód).} Przy utargu całkowitym $TR=P\cdot Q$ zachodzi kluczowa zależność z utargiem krańcowym $MR$:
\[
MR=\frac{dTR}{dQ}=P\!\left(1+\frac{1}{E_D}\right).
\]
Stąd praktyczna reguła cenowa:
\begin{itemize}
  \item popyt \emph{elastyczny} ($|E_D|>1$): obniżka ceny \emph{zwiększa} utarg ($MR>0$);
  \item popyt \emph{nieelastyczny} ($|E_D|<1$): podwyżka ceny \emph{zwiększa} utarg;
  \item $|E_D|=1$: utarg jest \emph{maksymalny} ($MR=0$).
\end{itemize}

\begin{definicja}[title={Inne elastyczności}]
\[
\underbrace{E_I=\frac{dQ}{dI}\frac{I}{Q}}_{\text{dochodowa}}\qquad
\underbrace{E_{XY}=\frac{dQ_X}{dP_Y}\frac{P_Y}{Q_X}}_{\text{mieszana (krzyżowa)}}\qquad
\underbrace{E_S=\frac{dQ_S}{dP}\frac{P}{Q_S}\ge 0}_{\text{cenowa podaży}}
\]
\begin{itemize}
  \item $E_I>0$ — dobro normalne; $E_I<0$ — dobro niższego rzędu (inferior); $E_I>1$ — dobro luksusowe.
  \item $E_{XY}>0$ — dobra substytucyjne; $E_{XY}<0$ — dobra komplementarne.
  \item $E_S$ rośnie z \emph{horyzontem czasu} (dłuższy okres $\Rightarrow$ większa elastyczność) i z wolnymi mocami produkcyjnymi.
\end{itemize}
\end{definicja}

\begin{uwaga}
Determinanty $|E_D|$: dostępność substytutów (więcej $\Rightarrow$ bardziej elastyczny), udział wydatku w budżecie, charakter dobra (dobra pierwszej potrzeby — nieelastyczne; luksusowe — elastyczne) oraz \emph{czas} na dostosowanie.
\end{uwaga}

% ---------------------------------------------------------------------
\subsection{Optimum techniczne i ekonomiczne przedsiębiorstwa}
% ---------------------------------------------------------------------

Rozważamy krótki okres z jednym czynnikiem zmiennym (praca $L$). Podstawowe krzywe produkcji:
\[
TP=f(L),\qquad AP=\frac{TP}{L}\ \text{(produkt przeciętny)},\qquad MP=\frac{dTP}{dL}\ \text{(produkt krańcowy)}.
\]
Od pewnego poziomu $L$ działa \textbf{prawo malejących przychodów krańcowych}: kolejne jednostki pracy dokładają coraz mniej produkcji ($MP$ maleje).

\begin{figure}[h]
\centering
\begin{tikzpicture}
\begin{groupplot}[
  group style={group size=1 by 2, vertical sep=1.15cm},
  width=11.5cm, height=4.3cm, xmin=0, xmax=8, domain=0:8, samples=120,
  axis lines=left, axis line style={-Stealth},
  every axis plot/.append style={very thick, smooth},
  xtick=\empty, ytick=\empty]
%
\nextgroupplot[ylabel={$TP$}, ymin=0, ymax=34, xlabel={}]
  \addplot[accent] {2*x^2 - 0.2*x^3};
  \node[accent] at (axis cs:7.2,20) {$TP$};
  \draw[dashed,gray] (axis cs:5,0) -- (axis cs:5,50);
  \draw[dashed,gray] (axis cs:6.67,0) -- (axis cs:6.67,50);
  \fill[accent] (axis cs:6.67,29.6) circle (2pt);
  \node[accent,above] at (axis cs:6.67,29.6) {\scriptsize $TP_{\max}$};
%
\nextgroupplot[ylabel={$AP,\;MP$}, ymin=0, ymax=8,
               xlabel={$L$ — nakład czynnika zmiennego}]
  \addplot[accentB] {2*x - 0.2*x^2};      \node[accentB] at (axis cs:7.4,3.0) {$AP$};
  \addplot[warn]    {4*x - 0.6*x^2};       \node[warn]    at (axis cs:2.2,6.55) {$MP$};
  \draw[dashed,gray] (axis cs:5,0) -- (axis cs:5,8);
  \draw[dashed,gray] (axis cs:6.67,0) -- (axis cs:6.67,8);
  \fill[accent] (axis cs:5,5) circle (2pt);
  \node[accent, font=\scriptsize, above left] at (axis cs:5,5) {$AP=MP$};
  \node[accent, font=\scriptsize, align=center, anchor=south] at (axis cs:4.55,5.9) {optimum\\ techniczne};
  \node[warn, font=\scriptsize, anchor=south] at (axis cs:6.67,6.2) {$MP=0$};
\end{groupplot}
\end{tikzpicture}
\caption{Krzywe produktu: całkowitego, przeciętnego i krańcowego. Optimum techniczne w~maksimum $AP$ (gdzie $AP=MP$).}
\label{fig:tpapmp}
\end{figure}

\begin{definicja}[title={Optimum techniczne vs ekonomiczne}]
\begin{itemize}
  \item \textbf{Optimum techniczne} — nakład, przy którym produkt przeciętny czynnika zmiennego jest \emph{największy} ($AP_{\max}$, gdzie $AP=MP$). To najbardziej \emph{wydajne technologicznie} wykorzystanie czynnika — zależy \emph{tylko} od technologii, nie od cen.
  \item \textbf{Optimum ekonomiczne} — nakład/produkcja \emph{maksymalizująca zysk}. Warunek: wartość produktu krańcowego = cena czynnika,
  \[
  MP_L\cdot P = w \quad\Longleftrightarrow\quad P = MC,
  \]
  gdzie $w$ — płaca, $P$ — cena produktu. Zależy \emph{od cen}, nie tylko od technologii.
\end{itemize}
\end{definicja}

Racjonalny producent działa w \emph{II strefie} produkcji — między optimum technicznym ($AP_{\max}$) a maksimum produktu całkowitego ($MP=0$). Optimum ekonomiczne leży wewnątrz tej strefy, a jego dokładne położenie wyznacza relacja cen $w/P$.

\begin{kluczowe}[title={Długi okres: reguła najniższego kosztu}]
Gdy zmienne są dwa czynniki (praca $L$, kapitał $K$), producent minimalizuje koszt danej produkcji tam, gdzie \emph{izokwanta} jest styczna do \emph{izokosztu}:
\[
MRTS=\frac{MP_L}{MP_K}=\frac{w}{r}\qquad\Longleftrightarrow\qquad
\frac{MP_L}{w}=\frac{MP_K}{r}.
\]
Czyli: ostatnia złotówka wydana na każdy czynnik daje ten sam przyrost produkcji.
\end{kluczowe}

% ---------------------------------------------------------------------
\subsection{Funkcja użyteczności i optimum konsumenta}
% ---------------------------------------------------------------------

\textbf{Funkcja użyteczności} $U(x,y)$ przypisuje koszykom dóbr liczbę tak, by lepszym koszykom odpowiadały większe wartości. Użyteczność jest \emph{porządkowa} (ordynalna) — liczy się \emph{ranking} koszyków, nie bezwzględne wartości.

\begin{definicja}[title={Użyteczność krańcowa i MRS}]
\[
MU_x=\frac{\partial U}{\partial x},\qquad
MRS_{xy}=\frac{MU_x}{MU_y}=-\left.\frac{dy}{dx}\right|_{U=\text{const}}.
\]
$MU$ zwykle \emph{maleje} (prawo malejącej użyteczności krańcowej). $MRS$ to nachylenie \emph{krzywej obojętności} — ile $y$ konsument odda za jednostkę $x$ przy stałej użyteczności.
\end{definicja}

Przy ograniczeniu budżetowym $P_x x + P_y y = I$ (dochód $I$) problem konsumenta to maksymalizacja $U$ na linii budżetu. Rozwiązanie leży w punkcie \emph{styczności} krzywej obojętności i linii budżetu (rys.~\ref{fig:konsument}):

\begin{kluczowe}[title={Warunek optimum konsumenta}]
\[
MRS_{xy}=\frac{MU_x}{MU_y}=\frac{P_x}{P_y}
\qquad\Longleftrightarrow\qquad
\boxed{\;\frac{MU_x}{P_x}=\frac{MU_y}{P_y}\;}
\]
(\emph{zasada wyrównywania użyteczności krańcowych na złotówkę}). Z mnożnikiem Lagrange'a $\mathcal L=U+\lambda(I-P_x x-P_y y)$ otrzymujemy $MU_i=\lambda P_i$, gdzie $\lambda$ to krańcowa użyteczność dochodu.
\end{kluczowe}

\begin{intuicja}
To po prostu warunek KKT/Lagrange'a: budżet przesuwamy między dobra, dopóki „przyrost zadowolenia na 1~zł'' nie wyrówna się we wszystkich dobrach. Gdyby był różny, przeniesienie złotówki do dobra o wyższym $MU/P$ podniosłoby użyteczność.
\end{intuicja}

\begin{figure}[h]
\centering
\begin{tikzpicture}[scale=1.0]
  \draw[->] (0,0) -- (6.2,0) node[right]{\small $x$ (dobro X)};
  \draw[->] (0,0) -- (0,6.2) node[above]{\small $y$ (dobro Y)};
  % linia budzetu
  \draw[very thick, accent] (0,5) -- (5,0);
  \node[accent, left]  at (0,5) {\small $\tfrac{I}{P_y}$};
  \node[accent, below] at (5,0) {\small $\tfrac{I}{P_x}$};
  % krzywe obojetnosci U = x*y
  \draw[very thick, accentB, domain=1.3:4.85, samples=100, smooth]
        plot (\x,{6.25/\x}) node[right]{\small $U^\ast$};
  \draw[accentB!55, domain=0.85:4.55, samples=100, smooth]
        plot (\x,{4/\x}) node[right]{\small $U_1$};
  % optimum
  \fill[accent] (2.5,2.5) circle (2.2pt) node[above right]{\small $E$};
  \draw[dashed, gray] (2.5,0) node[below]{\small $x^\ast$} -- (2.5,2.5) -- (0,2.5) node[left]{\small $y^\ast$};
  \node[accent, align=center] at (4.4,4.2) {\scriptsize styczność:\\[-2pt]\scriptsize $MRS=P_x/P_y$};
\end{tikzpicture}
\caption{Optimum konsumenta: styczność krzywej obojętności i linii budżetu.}
\label{fig:konsument}
\end{figure}

% ---------------------------------------------------------------------
\subsection{Maksymalizacja zysku i podaż firmy w konkurencji doskonałej}
% ---------------------------------------------------------------------

\begin{definicja}[title={Konkurencja doskonała — założenia}]
Wielu drobnych sprzedawców i nabywców; \emph{jednorodny} produkt; pełna informacja; swoboda wejścia i wyjścia; uczestnicy są \textbf{cenobiorcami} (price-takers). Firma sprzedaje po cenie rynkowej $P$, więc krzywa popytu na jej produkt jest \emph{pozioma}:
\[
P = AR = MR \quad(\text{cena} = \text{utarg przeciętny} = \text{utarg krańcowy}).
\]
\end{definicja}

Zysk to $\pi=TR-TC=P\,Q-TC(Q)$. Warunek maksymalizacji:
\[
\frac{d\pi}{dQ}=0\;\Rightarrow\; MR=MC \;\Rightarrow\; \boxed{\,P=MC(Q)\,},
\qquad\text{przy rosnącym } MC\ \Big(\tfrac{dMC}{dQ}>0\Big).
\]

\paragraph{Krótki okres.} O tym, czy firma produkuje, decyduje pokrycie kosztów \emph{zmiennych}:
\begin{itemize}
  \item produkuj, jeśli $P\ge \min AVC$ — inaczej \textbf{punkt zamknięcia} (shutdown): lepiej $Q=0$ i strata równa kosztowi stałemu;
  \item $P=\min ATC$ — \textbf{próg rentowności} (zysk zerowy);
  \item $P>ATC$ — zysk dodatni; $\min AVC< P<\min ATC$ — strata, ale firma produkuje (pokrywa część kosztu stałego).
\end{itemize}

\begin{kluczowe}[title={Krzywa podaży firmy (krótki okres)}]
To \emph{rosnąca część krzywej $MC$ powyżej $\min AVC$}. Poniżej $\min AVC$ podaż wynosi $0$. Podaż rynkowa = pozioma suma podaży firm.
\end{kluczowe}

\begin{figure}[h]
\centering
\begin{tikzpicture}[scale=1.0]
  \draw[->] (0,0) -- (8,0) node[right]{\small $Q$};
  \draw[->] (0,0) -- (0,6.4) node[above]{\small koszt, cena};
  % MC
  \draw[very thick, accent, smooth] plot coordinates
    {(0.5,4.6) (1,2.6) (1.8,1.4) (2.8,1.8) (3.6,2.8) (4.6,4.4) (5.3,6.0)};
  \node[accent] at (5.35,6.2) {\small $MC$};
  % AVC
  \draw[thick, warn, smooth] plot coordinates
    {(1,3.3) (1.8,2.2) (2.8,1.8) (3.8,2.15) (4.8,3.0) (5.4,3.8)};
  \node[warn] at (5.6,3.8) {\small $AVC$};
  % ATC
  \draw[thick, accentB, smooth] plot coordinates
    {(1.6,4.9) (2.4,3.6) (3.6,2.8) (4.4,3.05) (5.1,3.8) (5.5,4.3)};
  \node[accentB] at (5.75,4.35) {\small $ATC$};
  % cena
  \draw[very thick, accent!60!black, dashed] (0,4.55) -- (7.4,4.55);
  \node[accent!60!black, right] at (7.4,4.55) {\small $P=MR$};
  % Q*
  \fill[accent] (4.68,4.55) circle (2.2pt);
  \draw[dotted] (4.68,0) node[below]{\small $Q^\ast$} -- (4.68,4.55);
  % zysk (prostokat P .. ATC nad [0,Q*]) -- ATC(Q*)~3.0
  \draw[dotted] (0,3.0) -- (4.68,3.0);
  \fill[accent!12] (0,3.0) rectangle (4.68,4.55);
  \node[accent] at (2.3,3.8) {\small zysk};
  % punkty charakterystyczne
  \fill[warn]  (2.8,1.8) circle (2pt);   \node[warn, below] at (2.55,1.65) {\scriptsize $\min AVC$};
  \fill[accentB](3.6,2.8) circle (2pt);  \node[accentB, above left] at (3.55,2.9) {\scriptsize $\min ATC$};
\end{tikzpicture}
\caption{Równowaga firmy w krótkim okresie: $P=MC$, zysk = pole $(P-ATC)\times Q^\ast$. Podaż = $MC$ powyżej $\min AVC$.}
\label{fig:firma-ko}
\end{figure}

\begin{table}[h]
\centering
\small
\begin{tabular}{@{}ll@{}}
\toprule
\textbf{Relacja ceny} & \textbf{Sytuacja firmy (krótki okres)} \\
\midrule
$P>\min ATC$              & zysk ekonomiczny dodatni \\
$P=\min ATC$              & próg rentowności, zysk zerowy \\
$\min AVC<P<\min ATC$     & strata, ale produkuje \\
$P=\min AVC$              & punkt zamknięcia \\
$P<\min AVC$              & firma zamyka produkcję ($Q=0$) \\
\bottomrule
\end{tabular}
\caption{Decyzja produkcyjna firmy w~krótkim okresie.}
\end{table}

\paragraph{Długi okres.} Wszystkie czynniki są zmienne (koszt $LRAC$), a swoboda \emph{wejścia i wyjścia} zeruje zysk ekonomiczny: dodatni zysk przyciąga nowe firmy (podaż rośnie, cena spada), strata wypycha firmy (podaż maleje, cena rośnie). W równowadze długookresowej:
\[
P = MR = MC = \min LRAC = \min ATC.
\]

\begin{kluczowe}[title={Równowaga długookresowa}]
Firma produkuje w \emph{minimum} kosztu przeciętnego (efektywna skala), osiągając \textbf{zysk zerowy} (tzw. zysk normalny — koszt alternatywny kapitału jest już w kosztach). Dla gałęzi o stałych kosztach długookresowa krzywa podaży gałęzi jest \emph{pozioma} na poziomie $P=\min LRAC$.
\end{kluczowe}

\begin{figure}[h]
\centering
\begin{tikzpicture}[scale=0.95]
  \draw[->] (0,0) -- (7,0) node[right]{\small $Q$};
  \draw[->] (0,0) -- (0,5.6) node[above]{\small koszt, cena};
  % LRAC U-shaped
  \draw[very thick, accentB, smooth] plot coordinates
    {(0.7,4.6) (1.6,3.2) (3.0,2.5) (4.4,3.2) (5.4,4.6)};
  \node[accentB] at (5.6,4.7) {\small $LRAC$};
  % LMC through min LRAC
  \draw[thick, accent, smooth] plot coordinates
    {(1.2,1.3) (2.2,1.9) (3.0,2.5) (3.8,3.6) (4.6,5.1)};
  \node[accent] at (4.75,5.15) {\small $LMC$};
  % cena styczna do min
  \draw[very thick, accent!60!black, dashed] (0,2.5) -- (6.4,2.5);
  \node[accent!60!black, right] at (6.4,2.5) {\small $P=\min LRAC$};
  \fill[accent] (3.0,2.5) circle (2.2pt);
  \draw[dotted] (3.0,0) node[below]{\small $Q_{LR}$} -- (3.0,2.5);
  \node[accent, align=center] at (1.7,1.6) {\scriptsize zysk\\[-2pt]\scriptsize zerowy};
\end{tikzpicture}
\caption{Równowaga długookresowa: cena styczna do minimum $LRAC$; zysk ekonomiczny = 0.}
\label{fig:firma-do}
\end{figure}


\section{Polityka regulacji cen; wpływ opodatkowania na sytuację rynkową; dobra publiczne i dobra prywatne}
\textit{Treść w przygotowaniu.}

\section{Niedoskonałości rynku; asymetria informacyjna -- konsekwencje i metody redukcji}
\textit{Treść w przygotowaniu.}

\section{Formy konkurencji rynkowej -- konkurencja doskonała i konkurencja monopolistyczna, cechy rynku oligopolistycznego}
\textit{Treść w przygotowaniu.}

\section{Determinanty i konsekwencje nierówności dochodowych, krzywa Lorenza i współczynnik Giniego}
\textit{Treść w przygotowaniu.}

\section{Definicja i metody pomiaru PKB; PKB a PNB; ograniczenia PKB jako miernika dobrobytu; inflacja -- pomiar i przyczyny występowania; inflacja a deflator PKB; determinanty wzrostu gospodarczego w długim okresie; konwergencja; modele wzrostu gospodarczego}
\textit{Treść w przygotowaniu.}

\section{Bezrobocie -- definicja, rodzaje; pojęcie, rodzaje i funkcje pieniądza; mechanizm kreacji pieniądza}
\textit{Treść w przygotowaniu.}

\section{Międzynarodowe przepływy dóbr i kapitału oraz ich determinanty; zalety i wady teorii parytetu siły nabywczej i parytetu stóp procentowych; bilans handlowy z zagranicą; efekt mnożnikowy i efekt wypychania w keynesowskim modelu gospodarki; globalne łańcuchy wartości}
\textit{Treść w przygotowaniu.}

\section{Przyczyny, przebieg i skutki kryzysów gospodarczych; krzywa zagregowanego popytu oraz długo- i krótkookresowa krzywa zagregowanej podaży; krótko- i długookresowa krzywa Phillipsa; formy stabilizacyjnej polityki państwa}
\textit{Treść w przygotowaniu.}

\section{Modele przepływów międzygałęziowych; mnożniki input-output; globalne modele IO -- metoda konstrukcji, zalety, ograniczenia, zastosowania}
\textit{Treść w przygotowaniu.}

% =====================================================================
\chapter{Finanse}
% =====================================================================

\section{Rozwój światowych finansów i wycena aktywów w świetle historycznie postrzeganego ryzyka i niepewności; model równowagi ogólnej (Arrow (1951), Debreu (1951), Arrow, Debreu (1954))}
\textit{Treść w przygotowaniu.}

\section{Model Markowitza (1952) jako podwaliny teorii CAPM (1960--1962) i ICAPM Mertona (1973) -- krytyka, a warunki zgodności wyceny CAPM, ICAPM}
\textit{Treść w przygotowaniu.}

\section{ICAPM w świetle APT -- różnice i podobieństwa; alternatywne modele wyceny}
\textit{Treść w przygotowaniu.}

\section{Szacowanie ryzyka inwestycji w walory notowane na rynku giełdowym}
\textit{Treść w przygotowaniu.}

\section{Fundusze inwestycyjne -- ryzyko inwestycji w wybrane rodzaje funduszy}
\textit{Treść w przygotowaniu.}

\section{Możliwości zabezpieczenia działalności firmy w warunkach zmieniających się przyszłych stanów ekonomii}
\textit{Treść w przygotowaniu.}

\section{Ocena wyników sprawozdań finansowych spółek publicznych w świetle inwestora giełdowego}
\textit{Treść w przygotowaniu.}

\section{Wycena kapitałowa nieruchomości, przedsiębiorstwa -- metody i warunki dyskretyzacji}
\textit{Treść w przygotowaniu.}

\section{Rynkowa wartość obligacji, stopy zwrotu z inwestycji w obligacje w zależności od zmian rynkowych stóp procentowych}
\textit{Treść w przygotowaniu.}

\section{Zabezpieczenia wartości portfela obligacji -- duration obligacji, ryzyko okresu posiadania, ryzyko reinwestycji}
\textit{Treść w przygotowaniu.}

\end{document}
